%==============================================================================
% Chapitre 12 - Vent
%==============================================================================

\section{Introduction}

Le vent constitue l'un des paramètres environnementaux les plus importants en
balistique extérieure.

Son influence peut devenir prépondérante lorsque la distance augmente.

Contrairement à la gravité, dont les effets sont généralement prévisibles,
le vent varie :

\begin{itemize}
\item dans le temps ;
\item dans l'espace ;
\item en direction ;
\item en intensité.
\end{itemize}

Cette variabilité en fait l'une des principales sources d'incertitude lors
des tirs à moyenne et longue distance.

\begin{aretenir}

À longue distance, l'estimation du vent est souvent plus critique que la
connaissance de nombreux autres paramètres balistiques.

\end{aretenir}

\section{Qu'est-ce que le vent ?}

Le vent correspond à un déplacement de masses d'air dans l'atmosphère.

Il est généralement caractérisé par :

\begin{itemize}
\item sa vitesse ;
\item sa direction.
\end{itemize}

Dans ce manuel, la vitesse du vent sera généralement exprimée en m/s.

\section{Vent réel et vent apparent}

Le vent observé dépend du référentiel choisi.

Considérons un véhicule se déplaçant à 10~km/h dans un air parfaitement calme.

Pour un observateur immobile au sol :

\begin{itemize}
\item il n'y a pas de vent.
\end{itemize}

Pour un passager du véhicule :

\begin{itemize}
\item un vent apparent de 10~km/h semble souffler vers l'arrière.
\end{itemize}

Le projectile est sensible à la vitesse relative entre lui-même et l'air.

Cette notion est fondamentale en balistique.

\begin{exemple}

Un projectile quittant un véhicule en mouvement hérite initialement de la
vitesse du véhicule.

Cette vitesse doit être prise en compte dans le calcul de sa vitesse
relative par rapport à l'air.

\end{exemple}

\section{Vent de face}

Un vent de face souffle dans la direction opposée au mouvement du
projectile.

Il augmente généralement :

\begin{itemize}
\item la vitesse relative entre le projectile et l'air ;
\item la traînée aérodynamique ;
\item le temps de vol.
\end{itemize}

Le projectile ralentit davantage.

\section{Vent arrière}

Un vent arrière souffle dans la même direction que le projectile.

Il tend généralement à :

\begin{itemize}
\item diminuer la vitesse relative ;
\item réduire la traînée ;
\item améliorer légèrement la conservation de vitesse.
\end{itemize}

Son influence est souvent moins spectaculaire que celle d'un vent de
travers.

\section{Vent de travers}

Le vent de travers agit perpendiculairement à la trajectoire.

C'est généralement lui qui produit les erreurs de tir les plus visibles.

Sous son influence :

\begin{itemize}
\item le projectile dérive ;
\item le point d'impact se déplace latéralement.
\end{itemize}

Cette dérive augmente généralement avec :

\begin{itemize}
\item la vitesse du vent ;
\item le temps de vol ;
\item la distance de tir.
\end{itemize}

\begin{aretenir}

Le vent de travers est généralement le type de vent le plus critique pour
les tirs de précision.

\end{aretenir}

\section{Le vent agit-il directement sur la balle ?}

Une idée intuitive consiste à imaginer que le vent pousse le projectile
comme il pousserait une voile.

La réalité est plus complexe.

Le projectile interagit en permanence avec l'air qui l'entoure.

La trajectoire résulte de cette interaction aérodynamique.

La dérive observée est donc la conséquence de phénomènes aérodynamiques et
non d'une simple poussée latérale.

\section{Vent constant et vent variable}

Les modèles les plus simples supposent un vent constant sur toute la
trajectoire.

Dans la réalité, le vent varie souvent :

\begin{itemize}
\item avec l'altitude ;
\item avec le relief ;
\item avec le temps.
\end{itemize}

Ces variations compliquent considérablement la prévision des trajectoires.

\section{Gradient de vent}

La vitesse du vent n'est généralement pas identique à toutes les altitudes.

Cette variation est appelée gradient de vent.

À proximité du sol :

\begin{itemize}
\item les frottements ralentissent l'air ;
\item le vent est souvent plus faible.
\end{itemize}

En altitude :

\begin{itemize}
\item les frottements diminuent ;
\item le vent peut devenir plus rapide.
\end{itemize}

Un projectile peut donc traverser plusieurs couches d'air présentant des
conditions différentes.

\section{Lecture du vent}

L'estimation du vent constitue une compétence importante pour les tireurs.

Plusieurs indices peuvent être utilisés :

\begin{itemize}
\item végétation ;
\item fumée ;
\item poussière ;
\item manches à air ;
\item instruments météorologiques.
\end{itemize}

L'évaluation du vent reste toutefois une source d'incertitude importante.

\section{Influence sur les essais}

Les essais balistiques doivent systématiquement documenter :

\begin{itemize}
\item la vitesse du vent ;
\item sa direction ;
\item les éventuelles variations observées.
\end{itemize}

Ces informations sont essentielles pour interpréter correctement les
résultats.

\begin{simulation}

Les moteurs balistiques les plus simples utilisent un vent uniforme.

Les modèles avancés permettent de représenter :

\begin{itemize}
\item des couches atmosphériques ;
\item des gradients de vent ;
\item des vents variables dans le temps.
\end{itemize}

Ces modèles sont particulièrement utiles pour les simulations à longue
distance.

\end{simulation}

\section{Vers la balistique extérieure}

Le vent constitue l'un des principaux phénomènes étudiés en balistique
extérieure.

Les chapitres consacrés à la traînée et aux forces aérodynamiques
permettront de comprendre plus précisément les mécanismes responsables de la
dérive des projectiles.

\section{À retenir}

\begin{aretenir}

\begin{itemize}
\item Le vent est un déplacement de masses d'air.
\item Le projectile est sensible à sa vitesse relative par rapport à
l'air.
\item Le vent de travers est généralement le plus pénalisant.
\item Le vent réel varie souvent dans l'espace et dans le temps.
\item L'estimation du vent constitue une source majeure d'incertitude.
\end{itemize}

\end{aretenir}

